% GENERATED FILE. Edit canonical lessons, then run npm run generate:books.
% canonical-source-hash: fnv1a64:ebead3cbcf3714b4
% canonical-lessons: GU-C44-words, GU-C44-lines, GU-C44-pehli-vanchan

\chapter{Reading — Six Words, Six Lines, and a Whole Meeting}
\label{ch:gu-pehli-vanchan}
\hlchaptermodality{44}

\begin{chapteropening}
\textbf{By the end of this chapter:} \emph{I can read Gujarati off a signboard, read six lines I have only ever said aloud, and follow a whole meeting on the page without translating it word by word.}

The greetings, the repair and the parting the book taught first, read in the order two people would actually use them, with no new word in them.
\end{chapteropening}

\section[ghar · bajār · śāḷā · dukān · pāṇī …]{(reading) — six words off a signboard}

\label{lesson:GU-C44-words}



\begin{quote}
You have written every one of these. Not one is new.

What is new is that nobody is going to say them first.
\end{quote}

\begin{tcolorbox}[breakable,skin=enhanced,colback=orange!6,colframe=orange!45!black,boxrule=0.5pt,arc=1mm,left=6pt,right=6pt,top=4pt,bottom=4pt,fonttitle=\bfseries,title={Read this}]
\begin{quote}
\gu{ઘર} \gu{બજાર} \gu{શાળા} \gu{દુકાન} \gu{પાણી} \gu{પુસ્તક}
\end{quote}

Read down the list once. Do not build each one sign by sign — look at the whole word and let it arrive.

Again, and notice which came at once and which you had to assemble. Both are fine. The difference is what practice removes.
\end{tcolorbox}

\subsection*{What to know first}
Writing gave you the signs in order and all the time you wanted. A signboard gives you one look.

These six are the ones you would actually meet painted on something — over a shop, on a school gate, beside a tap.

That is why six is enough for one sitting.

\subsection*{Your turn}
\begin{itemize}
\raggedright
  \item \emph{Say it:} the six words down the list, once, without stopping
\end{itemize}

\emph{Twice through:}

\begin{tcolorbox}[breakable,colback=teal!4,colframe=teal!35!black,title={Before you move on}]
How many of the six were new? (\textbf{None.})

What was new? (\textbf{Nobody said them first.} You read them.)
\end{tcolorbox}
\section[mārũ nām mīrā chhe. tame kem chho?]{(reading) — six lines, no voice attached}

\label{lesson:GU-C44-lines}



\begin{quote}
The last lesson gave you single words. These are whole lines.

You have said every one of them out loud. None of them is new.
\end{quote}

\begin{tcolorbox}[breakable,skin=enhanced,colback=orange!6,colframe=orange!45!black,boxrule=0.5pt,arc=1mm,left=6pt,right=6pt,top=4pt,bottom=4pt,fonttitle=\bfseries,title={Read this}]
\begin{quote}
\gu{મારું} \gu{નામ} \gu{મીરા} \gu{છે।} \gu{તમે} \gu{કેમ} \gu{છો}? \gu{હું} \gu{મજામાં} \gu{છું।} \gu{વાંધો} \gu{નહીં।} \gu{માફ} \gu{કરો।} \gu{કાલે} \gu{મળીશું।}
\end{quote}

Read down once, without stopping.

Again — and this time notice that nobody said them first. Until now every one of these arrived in your ear before it arrived on the page.
\end{tcolorbox}

\subsection*{What to know first}
These six are not a list. Read in order they are most of a conversation: a name, a question, an answer, a repair, an apology, a parting.

That is worth seeing once on the page. Spoken, each of them came with a voice telling you what it was for. Written, the order has to do that work instead.

Six lines, and the order is the whole lesson.

\subsection*{Your turn}
\begin{itemize}
\raggedright
  \item \emph{Say it:} the six lines, once, in order, without stopping
\end{itemize}

\emph{Twice through:}

\begin{tcolorbox}[breakable,colback=teal!4,colframe=teal!35!black,title={Before you move on}]
How many of the six lines were new? (\textbf{None.})

What did reading them in order show you? (\textbf{They are a conversation}, not a list.)
\end{tcolorbox}
\section[tame kem chho? hũ majāmā̃ chhũ.]{(reading) — a whole meeting}

\label{lesson:GU-C44-pehli-vanchan}



\begin{quote}
Six words, then six lines. Now two people, from hello to goodbye.

Thirty-five words, and not one of them is new.
\end{quote}

\begin{tcolorbox}[breakable,skin=enhanced,colback=orange!6,colframe=orange!45!black,boxrule=0.5pt,arc=1mm,left=6pt,right=6pt,top=4pt,bottom=4pt,fonttitle=\bfseries,title={Read this}]
\begin{quote}
\gu{તમે} \gu{કેમ} \gu{છો}? \gu{હું} \gu{મજામાં} \gu{છું।} \gu{તમારું} \gu{નામ} \gu{શું} \gu{છે}? \gu{મારું} \gu{નામ} \gu{મીરા} \gu{છે।} \gu{તમને} \gu{મળીને} \gu{આનંદ} \gu{થયો।} \gu{આ} \gu{મારો} \gu{ભાઈ} \gu{છે।} \gu{આ} \gu{મારું} \gu{ઘર} \gu{છે।} \gu{માફ} \gu{કરો}, \gu{હું} \gu{સમજતો} \gu{નથી।} \gu{વાંધો} \gu{નહીં।} \gu{કાલે} \gu{મળીશું।}
\end{quote}

Read it through once without stopping. Do not translate as you go — let the turns arrive.

Now again, and count how many people are speaking.
\end{tcolorbox}

\subsection*{What to know first}
Two. Nothing on the page says so — no names in front of the lines, no he-said. The turns alternate and you followed them.

Notice the eighth line. \textbf{\gu{માફ} \gu{કરો}, \gu{હું} \gu{સમજતો} \gu{નથી।}} One of them has failed to understand, and the conversation carries on anyway. That is why the book taught the repair before it taught much else.

This is the first whole exchange you have read rather than said, and it works because every line was paid for, one chapter at a time.

\subsection*{Your turn}
\begin{itemize}
\raggedright
  \item \emph{Say it:} the whole exchange aloud, once, without stopping
\end{itemize}

\emph{Twice through:}

\begin{tcolorbox}[breakable,colback=teal!4,colframe=teal!35!black,title={Before you move on}]
How many words in that exchange were new? (\textbf{None.})

What told you where one speaker stopped? (\textbf{The order} — nothing else did.)
\end{tcolorbox}
